% paper/oe_supplement.tex -- Phase 7.4 supplement.
% Same template-pending status as oe_main.tex (see that file's header).
% Uses numbers.tex macros throughout, same as the main manuscript.

\documentclass[10pt]{article}
\usepackage[margin=1in]{geometry}
\usepackage{amsmath,amssymb,graphicx,booktabs,hyperref}
% AUTO-GENERATED by scripts/make_numbers_tex.py -- DO NOT HAND-EDIT.
% Regenerate: python scripts/make_numbers_tex.py
\usepackage{xcolor}  % for the [PENDING] flag color, harmless if already loaded

\newcommand{\MeasuredContrastTwoX}{2.03}
\newcommand{\KcPredTwoX}{4.25}
\newcommand{\KstarTwoXInterp}{\textcolor{red}{[PENDING]}}
\newcommand{\KstarTwoXCI}{7.85}
\newcommand{\MeanGainTwoX}{0.34}
\newcommand{\MeasuredContrastFourX}{3.98}
\newcommand{\KcPredFourX}{6.96}
\newcommand{\KstarFourXInterp}{\textcolor{red}{[PENDING]}}
\newcommand{\KstarFourXCI}{15.71}
\newcommand{\MeanGainFourX}{0.65}
\newcommand{\MeasuredContrastEightX}{6.03}
\newcommand{\KcPredEightX}{7.39}
\newcommand{\KstarEightXInterp}{\textcolor{red}{[PENDING]}}
\newcommand{\KstarEightXCI}{15.71}
\newcommand{\MeanGainEightX}{0.73}
\newcommand{\SeedCount}{\textcolor{red}{[PENDING]}}
\newcommand{\NResultFiles}{684}

% --- supplement macros (already-real data, not Phase-3-gated) ---
\newcommand{\AblationCheckpointCosSim}{0.977}
\newcommand{\AblationSurrogateCosSim}{0.841}
\newcommand{\RCWAMaxDevThreeCase}{0.038}
\newcommand{\RCWAE7MaxDev}{0.980}
\newcommand{\RCWAE7NCases}{90}
\newcommand{\MeshPSNRSpread}{0.036}
\newcommand{\WavelengthDesignPSNR}{6.63}


\title{Supplement: Media-in-the-Loop Holography}
\date{}

\begin{document}
\maketitle

\section{Gradient-Pathway Ablation}
Checkpointed discrete-adjoint gradients match unrolled-autodiff gradients
to cosine similarity \AblationCheckpointCosSim\ on a realistic
reconstruction-loss probe; a neural surrogate reaches
\AblationSurrogateCosSim. See \texttt{experiments/ablation\_gradients.py},
\texttt{results\_ablation\_gradients.json}, Fig.~F8a.

\section{RCWA Validity Envelope}
The original 3-case scalar-vs-vector cross-check (TE, unslanted,
near-Bragg incidence) agrees to \RCWAMaxDevThreeCase\ maximum absolute
diffraction-efficiency deviation. The E7 extension (TE+TM, 3
incidence/slant geometries, 3 $\Delta n$ levels, 5 $K$ values --
\RCWAE7NCases\ cases total) finds \RCWAE7MaxDev\ maximum deviation,
concentrated in the slanted-grating geometry at high $\Delta n$: a
20-degree slant detunes the Bragg condition at the nominal (unslanted-
Bragg) incidence angle, so RCWA correctly shows near-zero diffraction
while the \emph{unslanted} Kogelnik formula (no slant term) predicts
near-peak efficiency. This is a genuine finding about the unslanted
formula's validity envelope, not a numerical artifact -- see
\texttt{results\_rcwa\_e7.json}'s \texttt{by\_geometry} breakdown and
Fig.~F8b. Readout caveat: the compensation cliff (main text, Eq.~5) is a
\emph{recording}-side phenomenon; this grid characterizes
\emph{readout}-side scalar-vs-vector error, which moves absolute PSNRs
more than paired (M4$-$M2) gains -- E7 scopes the readout engine's
validity envelope, it does not itself validate or invalidate the
cliff-vs-budget result.

\section{GPU Mesh-Density Convergence}
Reconstruction PSNR is mesh-independent to within \MeshPSNRSpread~dB
across a $4\times$ resolution range ($n_x=512/1024/2048$, fixed 51.2\,$\mu$m
physical window); wall-clock was essentially flat across the same range on
a single Colab T4 (single seed) -- the optimizer is overhead-bound, not
compute-bound, at this problem size, which motivated the batched
implementation used for all Phase 3 manifests. See Fig.~F8c,
\texttt{results/gpu\_reruns/npdd\_mesh\_sweep/results.json}.

\section{Wavelength-Detuning Readout Sweep}
A hologram recorded at the 405\,nm design wavelength reaches
\WavelengthDesignPSNR~dB at that wavelength; reading the same recorded
profile out at 400--450\,nm shows a \emph{non-monotonic} falloff away
from the design point, consistent with Kogelnik's oscillatory
$\sin^2(\nu)$ angular/wavelength detuning response rather than a smooth
rolloff (single seed, single Colab T4 run). See Fig.~F8d,
\texttt{results/gpu\_reruns/bpm\_wavelength\_sweep/results.json}.

\section{Full Result Tables}
\textbf{[TODO, blocked on Phase 3: full E1-E6 result tables, pending
manifest runs. Regenerate this section from
\texttt{results/summary/paper\_numbers.json} once populated --
\texttt{analysis/aggregate.py}'s \texttt{per\_config} dict has every
(experiment, config, method) cell needed.]}

\section{Provenance Appendix}
Every number in the main text and this supplement traces to a committed
results JSON via \texttt{results/summary/paper\_numbers.json}
(\texttt{analysis/aggregate.py}) or, for the supplement-only quantities
above, directly to the named JSON file in each section. No number in
either document is hand-typed -- both are built exclusively from
\texttt{numbers.tex} (\texttt{scripts/make\_numbers\_tex.py}), and
\texttt{scripts/check\_consistency.py} fails the build on any violation
of that rule.

\end{document}
